\documentclass[12pt]{article}
\usepackage[utf8]{inputenc}
\usepackage[T1]{fontenc}
\usepackage{hyperref}
\title{Kaamba Language Resources for K-12 Curriculum Development}
\author{Compiled by Claude Sonnet 5.5 (claude-sonnet-5-5) for Andr\'as Kornai}
\date{}
\begin{document}
\maketitle

\section{Introduction}
Kaamba (ISO 639-3 xku; Glottocode kaam1238) is a Bantu language of the Kongo cluster, spoken in the Republic of the Congo (Brazzaville), in the Bouenza and L\'ekoumou regions. Sources give about 2,900 speakers (Global Recordings) and about 3,600 (Joshua Project). Do not confuse it with the Kamba of Kenya, a much larger and unrelated Northeast Bantu language. Search found very little material.

\section{Primary Educational Resources}
No primers, textbooks or school programmes for Kaamba were found in search results.

\section{Linguistic and Reference Materials}
\begin{itemize}
\item Wikipedia, ``Kongo languages'': background on the Kongo cluster to which Kaamba belongs. \url{https://en.wikipedia.org/wiki/Kongo_languages}
\end{itemize}

\section{Religious and Community Resources}
\begin{itemize}
\item Joshua Project, Kaamba in Congo: profile. It reports that Bible translation has started but no Bible, Jesus Film or Gospel recordings are available. \url{https://www.joshuaproject.net/people_groups/19175}
\item Joshua Project language page for xku. \url{https://joshuaproject.net/languages/xku}
\end{itemize}

\section{Digital and Audio Resources}
\begin{itemize}
\item Global Recordings Network page for xku. Global Recordings also lists Doondo under the name Kaamba, which may be a different language, so check before use. \url{https://globalrecordings.net/en/language/xku}
\end{itemize}

\section{Curriculum Development Context}
No evidence of schooling in Kaamba was found. French is the national language of the Republic of the Congo. A curriculum would require basic documentation and orthography work first, and could draw on the better-documented Kongo languages.
\end{document}